\documentclass[10pt,a4paper]{amsart}
\usepackage[margin=19mm]{geometry}
\usepackage{amsmath,amssymb,mathtools}
\usepackage[scr=rsfs,cal=euler]{mathalfa}
\usepackage{xfrac}
\usepackage[unicode,hidelinks,bookmarks=false]{hyperref}
\newcommand{\ie}{i.e.\ }
\newcommand{\cf}{cf.\ }
\newcommand{\resp}{respectively\ }
\newcommand{\Subsection}{\subsection}
\providecommand{\nobreakdash}{\nobreak\hbox{-}}
\setlength{\emergencystretch}{2em}
\multlinegap0pt
\usepackage{bm}
\usepackage{bbm}
\usepackage{dsfont}
\usepackage{xspace}
\usepackage{cancel}
\usepackage{mathtools}
\usepackage{cleveref}
\usepackage{amsthm}
\makeatletter
\@ifundefined{lemma}{\newtheorem{lemma}{Lemma}}{}
\makeatother
\newcommand{\CC}{\mathbb{C}}
\newcommand{\RR}{\mathbb{R}}
\newcommand{\T}{\mathbb{T}}
\newcommand{\abs}[1]{\left\lvert #1\right\rvert}
\newcommand{\paren}[1]{\left(#1\right)}
\title[Graphs with connectivity $3/4 - \varepsilon$ are globally sy]{Graphs with connectivity $3/4 - \varepsilon$ are globally synchronizing}
\author{Saba Lepsveridze, Sam Zhang}
\date{Source: arXiv:2608.20010v1 (CC BY 4.0)}
\begin{document}
\maketitle

\noindent\emph{Source and licence.} Graphs with connectivity $3/4 - \varepsilon$ are globally synchronizing. Version arXiv:2608.20010v1 (CC BY 4.0), distributed under the Creative Commons Attribution 4.0 International licence, as recorded in the arXiv licence metadata for this exact version and shown on the abstract page https://arxiv.org/abs/2608.20010v1. Full rights notice: licenses/arxiv-2608.20010-CC-BY-4.0.txt.\par\smallskip

\noindent\textit{Editorial scope.} Counted unit: the Lemma whose source label is \texttt{lem:exceptional-vertex-main} in the article (source0.tex lines 906--922, source label \texttt{lem:exceptional-vertex-main}), delivered together with the article's own complete original proof of it (source0.tex lines 924--998, one \texttt{proof} environment of 75 lines). No mathematics is written, repaired or re-derived: statement and proof are the article's own text line for line. Required context retained uncounted: lem:exceptional-ideal-case (lines 841--853). Every definition, display and label of those lines is retained unchanged. Omitted from this package and not redistributed: 1--840, 854--905, 923--923, 999--1125. Nothing inside a retained excerpt is omitted or altered: every retained body is delivered exactly as the article's lines are written, so no omission receipt of any kind accompanies this package. Citations: the retained excerpts carry no citation command at all, so no bibliography entry of the article is transcribed and no reference is redistributed. Reference adaptation: none was needed and none was made; every reference inside the delivered unit resolves inside it, so no unresolved reference or citation is produced. Printed slips kept literal: no printed word, symbol, hypothesis or number of the counted statement or of its proof was altered. Layout and class: the article's own class is \texttt{amsart} with packages including geometry, fontenc, amsmath, amssymb, mathtools, hyperref, cleveref, xcolor; the wrapper is the lane's pinned \texttt{amsart} editorial wrapper at 10pt on A4, which restates the theorem-like environments the retained text needs and adds the article's own macro definitions that the retained text needs (\texttt{\textbackslash CC}, \texttt{\textbackslash RR}, \texttt{\textbackslash T}, \texttt{\textbackslash abs}, \texttt{\textbackslash paren}), with extra wrapper packages (bm, bbm, dsfont, xspace, cancel, mathtools, cleveref). The delivered numbering is the unit's own. Assets: no retained span contains a figure environment, a graphics-inclusion command, a PDF-page inclusion, a table or a TikZ picture; no image file is copied into this package, the package declares no asset dependency and no image file is redistributed with it. Licence: version arXiv:2608.20010v1 of this article is distributed under the Creative Commons Attribution 4.0 International licence, as recorded in the arXiv licence metadata for this exact version and shown on the abstract page https://arxiv.org/abs/2608.20010v1. A full rights notice is delivered at licenses/arxiv-2608.20010-CC-BY-4.0.txt. Characters: every retained excerpt is pure ASCII, so the delivered engine, XeLaTeX with its default Latin Modern text font, reports no missing character.
\begin{lemma}\label{lem:exceptional-ideal-case}
For all $q_0 > 1/\sqrt2$ and $h > 0$, there is a constant $\kappa>0$, such that the following  holds.  Suppose $\phi\in\T$ is at least $h/2$ from every multiple of $\pi/2$, and $W_0,W_1,W_2,W_3\in[0,1/4]$ satisfy
\[
    \sum_{j=0}^3W_j\ge q_0,
    \qquad
    \sum_{j=0}^3W_j\sin\paren{\frac{\pi j}{2}-\phi}=0.
\]
Then there is a $z\in\CC$ with $\abs z=1$ such that
\[
    \sum_{j=0}^3W_j\cos\paren{\frac{\pi j}{2}-\phi}
       \abs{z-i^j}^2\le-\kappa.
\]
\end{lemma}

\begin{lemma}\label{lem:exceptional-vertex-main}
There is a constant $\tau_0>0$, depending only on $h$ and $q_0$, such that the
following holds.  Suppose $(G,\theta)$ is
$(h,\rho,\alpha,\beta)$-clustered with $\mu(G)\ge3/4-\eta$, and
\[
    0\le\eta,\rho,\alpha,\beta\le\tau_0.
\]
For every exceptional vertex $u\in E$, there is a $z_u\in\CC$ with
$\abs{z_u}=1$ such that
\begin{equation}\label{eq:exceptional-vertex-negative}
\sum_{j=0}^3\sum_{v\in B_j}
w_{uv}\cos(\theta_v-\theta_u)
\abs{z_u-i^j(1+it_v)}^2
\le-\frac{\kappa n}{2},
\end{equation}
where $\kappa$ is from \cref{lem:exceptional-ideal-case}.
\end{lemma}

\begin{proof}
Fix $u\in E$, put $\phi=\theta_u-\bar t$, and define
\[
    W_j:=\frac1n\sum_{v\in B_j}w_{uv}.
\]
Then $0\le W_j\le1/4+\rho$ and $W_0 + W_1 + W_2 + W_3 \geq 3/4 - \eta - \beta$. Since $\theta_v-\theta_u=\pi j/2-\phi+t_v$ for $v\in B_j$, the
first-order condition and the $L^1$ bound on the $t_v$ give
\begin{equation}\label{eq:perturbed-first-order}
\Bigl|{\sum_{j=0}^3W_j\sin\Bigl({\frac{\pi j}{2}-\phi}\Bigr)}\Bigr|
\le\alpha+\beta.
\end{equation}
We next perturb these weights to impose exact first-order balance.  For
$X\in\RR^4$, write
\[
    D_\phi(X):=\sum_{j=0}^3X_j
        \sin\paren{\frac{\pi j}{2}-\phi}.
\]
Set $\lambda=(1+4\rho)^{-1}$ and $W'_j=\lambda W_j$.  Then
$W'_j\in[0,1/4]$,
\[
    \sum_j\abs{W_j-W'_j}\le4\rho,
    \qquad
    D_\phi(W')=\lambda D_\phi(W).
\]
Put $s=\sin\phi$, $c=\cos\phi$, $d=W'_2-W'_0$, and
$e=W'_1-W'_3$, so that $D_\phi(W')=sd+ce$.  If
$\abs s\ge\abs c$, replace $(W'_0,W'_2)$ by a pair
$(\tilde W_0,\tilde W_2)\in[0,1/4]^2$ whose difference is
$-ce/s$, leaving the other two coordinates unchanged.  Such a pair exists
because $\abs{ce/s}\le1/4$, and it may be chosen at $\ell^1$ cost
\[
    \abs{(-ce/s)-d}
    =\frac{\abs{D_\phi(W')}}{\abs s}
    \le\sqrt2|{D_\phi(W')}|.
\]
If $\abs c>\abs s$, the symmetric construction instead adjusts
$(W'_1,W'_3)$.  In either case, we have
\[
    \tilde W_j\in[0,1/4],
    \qquad
    D_\phi(\tilde W)=0,
    \qquad
    \sum_j|{W_j-\tilde W_j}|
    \le K(\rho+\alpha+\beta)
\]
for an absolute constant $K$, where we used
\eqref{eq:perturbed-first-order} in the last inequality. For sufficiently small $\tau_0$, the perturbation estimate above imply
\[
    \sum_j\tilde W_j
    \ge\frac34-\eta-\beta-K(\rho+\alpha+\beta)\ge q_0.
\]
Moreover, because $u\in E$ and
$\abs{\bar t}\le h/2$ for small enough $\alpha$, the phase $\phi$ is at least $h/2$ from every
multiple of $\pi/2$.  We may therefore apply
\cref{lem:exceptional-ideal-case} and obtain $z_u$, with
$\abs{z_u}=1$, such that
\begin{equation}\label{eq:tilde-w-upper-bound}
\sum_j\tilde W_j\cos\paren{\frac{\pi j}{2}-\phi}
\abs{z_u-i^j}^2\le-\kappa.
\end{equation}

By the choice of $h_0$, we have $|{t_v}|\le h+\alpha<1$.  Because
$|{z_u-i^j}|\le2$, the difference between
$|{z_u-i^j(1+it_v)}|^2$ and $|{z_u-i^j}|^2$ is bounded by an absolute
constant times $\abs{t_v}$.
The cosine function is Lipschitz, and $|{(\theta_v-\theta_u)-(\pi j/2-\phi)}|=|{t_v}|.$ Consequently, using \eqref{eq:tilde-w-upper-bound},
\begin{equation*}
\frac1n\sum_j\sum_{v\in B_j}
w_{uv}\cos(\theta_v-\theta_u)
\abs{z_u-i^j(1+it_v)}^2\le-\kappa+C_2(\rho+\alpha+\beta)
\end{equation*}
for an absolute constant $C_2$.
After decreasing $\tau_0$, the right-hand side is at most
$-\kappa/2$.
\end{proof}

\end{document}
